\documentclass[11pt]{article}

\usepackage[utf8]{inputenc}
\usepackage[T1]{fontenc}
\usepackage[margin=1in]{geometry}
\usepackage{natbib}
\usepackage{hyperref}
\usepackage{url}
\usepackage{parskip}

\title{\texttt{mushin}: Boilerplate-Free, Reproducible\\
       Machine-Learning Experiment Sweeps}

% TODO: add affiliation and \thanks with your ORCID before submission.
\author{Josu\'e Mart\'{\i}nez-Mart\'{\i}nez}

\date{\today}

\begin{document}
\maketitle

\begin{abstract}
\texttt{mushin} is a Python library for running reproducible machine-learning
experiment sweeps with minimal boilerplate. A researcher decorates an experiment
function with \texttt{@mushin.sweep}, declares the parameters to sweep, and
receives the results as a labeled \texttt{xarray.Dataset}~\citep{hoyer2017xarray}
keyed by the swept dimensions --- rather than as rows in a dashboard that must be
exported and re-assembled. The sweep layer is framework-agnostic: a task returns
a dictionary of metrics, so scikit-learn~\citep{pedregosa2011scikit}, XGBoost, or
any Python model can be swept the same way as a PyTorch~\citep{paszke2019pytorch}
model. Built on PyTorch Lightning~\citep{falcon2019lightning} and
hydra-zen~\citep{hydrazen2022}, \texttt{mushin} records each run's configuration
and provenance, resumes durably after interruption, and scales from a laptop to a
multi-node SLURM cluster by changing only the launcher. It is available on PyPI
as \texttt{mushin-py}.
\end{abstract}

\section{Introduction}

Hyperparameter and seed sweeps are central to empirical machine-learning
research, but the surrounding infrastructure is repetitive and error-prone.
Practitioners commonly stitch together a configuration system, a launcher, a
results store, and analysis glue by hand. The results often end up in an
experiment-tracking dashboard~\citep{biewald2020wandb, mlflow2018} from which they
must be exported before analysis, and the exact configuration that produced a
given number is easy to lose.
Hyperparameter-optimization frameworks~\citep{akiba2019optuna, liaw2018tune} focus
on \emph{searching} for a single best configuration rather than producing the
full labeled grid of results that scientific comparison and ablation require, and
they still leave provenance, resumption, and multi-GPU scaling to the user.

\texttt{mushin} targets this gap with a small, composable workflow layer. Its core
idiom returns the Cartesian product of swept parameters directly as a labeled
dataset:

\begin{quote}
\begin{verbatim}
import mushin

@mushin.sweep
def experiment(lr, seed):
    ...                      # train and evaluate
    return dict(accuracy=acc)

ds = experiment.run(
    lr=mushin.multirun([0.01, 0.1, 1.0]),
    seed=mushin.multirun([0, 1, 2]),
)                            # xarray.Dataset, dims (lr: 3, seed: 3)
ds["accuracy"].mean("seed")  # reduce over seeds, per learning rate
\end{verbatim}
\end{quote}

\section{Design and features}

\texttt{mushin} provides four properties that are usually absent together in a
single tool.

\paragraph{Sweep $\rightarrow$ labeled dataset.}
The result of a sweep is an \texttt{xarray.Dataset}, so per-seed reduction,
slicing, and plotting are immediate and the mapping from configuration to result
is never lost. Because the task only returns a \texttt{dict}, the sweep interface
is framework-agnostic and results are standard \texttt{xarray} objects that
compose with the existing scientific-Python stack.

\paragraph{Reproducibility by construction.}
Each run's resolved configuration and provenance are captured automatically.
Auto-tuning helpers run PyTorch Lightning's tuner once and \emph{pin} the result
(e.g.\ a batch size, with gradient accumulation) to a sidecar file, so a run is
reproducible across hardware rather than silently re-tuning. Sweeps are durably
resumable: a sweep interrupted mid-run --- including by a hard process kill or a
real SLURM preemption --- resumes without recomputing completed cells.

\paragraph{Laptop-to-cluster with one code path.}
The same task runs in-process, across local worker processes (joblib), or across
SLURM nodes (submitit~\citep{submitit2021}) by changing only a launcher argument,
built on Hydra's launcher plugins~\citep{yadan2019hydra}. Data-parallel (DDP) and
fully-sharded (FSDP) multi-GPU training, and round-robin GPU packing for small
jobs, are supported and have been validated on real multi-node GPU hardware.

\paragraph{Statistics-aware comparison.}
Built-in benchmark ``batteries'' and a \texttt{compare} API evaluate methods
across seeds with significance testing, so a claimed improvement comes with a
measure of uncertainty rather than a single number.

\section{Related work}

Experiment trackers such as Weights \& Biases~\citep{biewald2020wandb} and
MLflow~\citep{mlflow2018} record runs but leave sweep orchestration and labeled
result assembly to the user. Hyperparameter optimizers such as
Optuna~\citep{akiba2019optuna} and Ray Tune~\citep{liaw2018tune} search for a best
configuration rather than returning the full labeled grid. Hydra~\citep{yadan2019hydra}
and hydra-zen~\citep{hydrazen2022} provide configuration and multirun machinery;
\texttt{mushin} builds on them to add the sweep-to-dataset abstraction,
reproducibility features, and validated multi-GPU/multi-node scaling.

\section{Availability}

\texttt{mushin} is open source under the MIT license, installable from PyPI
(\texttt{pip install mushin-py}), with documentation, runnable examples, and a
test suite at \url{https://github.com/martinez-hub/mushin}. Each release is
archived on Zenodo; the concept DOI
\href{https://doi.org/10.5281/zenodo.21436444}{10.5281/zenodo.21436444} always
resolves to the latest version~\citep{mushin_zenodo}.

\section*{Acknowledgements}

\texttt{mushin} is a maintained, standalone carve-out of the \texttt{mushin}
subpackage of MIT Lincoln Laboratory's
\texttt{responsible-ai-toolbox}~\citep{raitoolbox2021}; this project maintains and
extends that workflow layer, whose original design originates from that toolbox.

\bibliographystyle{plainnat}
\bibliography{paper}

\end{document}
